% GENERATED FILE. Edit canonical lessons, then run npm run generate:books.
% canonical-source-hash: fnv1a64:90a87eb74ead2bf4
% canonical-lessons: KA-C401-abhivrddhi, KA-C401-arthika, KA-C401-hanakasu, KA-C401-uttejisu, KA-C401-hudu

\chapter{Growth and Finance}
\label{ch:ka-growth-and-finance}
\hlchaptermodality{401}

\begin{chapteropening}
\textbf{By the end of this chapter:} \emph{I can say development, economic, finance, to promote and to invest.}

Complete the last lesson of chapter 401: I can say development, economic, finance, to promote and to invest.
\end{chapteropening}

\section[\texorpdfstring{\kn{ಅಭಿವೃದ್ಧಿ} (abhivṛddhi)}{ಅಭಿವೃದ್ಧಿ (abhivṛddhi)}]{\kn{ಅಭಿವೃದ್ಧಿ} (abhivṛddhi) — development}

\label{lesson:KA-C401-abhivrddhi}



\begin{quote}
Before the new one: say the Kannada for to utilise, then the Kannada for to waste.
\end{quote}

\subsection*{You'll want to know: \kn{ಅಭಿವೃದ್ಧಿ}}
\textbf{\kn{ಅಭಿವೃದ್ಧಿ}} — \emph{abhivṛddhi} — \textquotedblleft{}development\textquotedblright{}.

\begin{quote}
\kn{ಹೊಸ} \kn{ರಸ್ತೆಯಿಂದ} \kn{ಊರಿನ} \kn{ಅಭಿವೃದ್ಧಿ} \kn{ಆಯಿತು}.
\end{quote}

\begin{quote}
\emph{hosa rasteyinda ūrina abhivṛddhi āyitu.}
\end{quote}

\begin{quote}
The new road brought development to the town.
\end{quote}

\kn{ಅಭಿವೃದ್ಧಿ} \kn{ಹೊಂದು}, to develop.

\begin{cousinweb}
\textbf{\kn{ಅಭಿವೃದ್ಧಿ}} (\emph{abhivṛddhi}) is borrowed from Sanskrit \emph{abhivṛddhi}, \textquotedblleft{}growth, increase, prosperity\textquotedblright{} (Monier-Williams).
\end{cousinweb}

\begin{grammarlens}[title={What you've built}]
One more word for giving reasons and opinions.
\end{grammarlens}

\subsection*{Your turn}
\begin{itemize}
\raggedright
  \item \emph{Say it:} \emph{abhivṛddhi}
  \item \emph{Say it:} \emph{abhivṛddhi}, once more
  \item \emph{Say it:} \emph{hosa rasteyinda ūrina abhivṛddhi āyitu.} — The new road brought development to the town.
  \item \emph{Recall:} say the Kannada for to utilise, then the Kannada for to waste, then say \emph{abhivṛddhi} again
\end{itemize}

\begin{tcolorbox}[breakable,colback=teal!4,colframe=teal!35!black,title={Before you move on}]
What does \kn{ಅಭಿವೃದ್ಧಿ} mean? (\textquotedblleft{}Development\textquotedblright{}.) Say it once more.
\end{tcolorbox}
\section[\texorpdfstring{\kn{ಆರ್ಥಿಕ} (ārthika)}{ಆರ್ಥಿಕ (ārthika)}]{\kn{ಆರ್ಥಿಕ} (ārthika) — economic; financial}

\label{lesson:KA-C401-arthika}



\begin{quote}
Before the new one: say the Kannada for to waste, then the Kannada for development.
\end{quote}

\subsection*{You'll want to know: \kn{ಆರ್ಥಿಕ}}
\textbf{\kn{ಆರ್ಥಿಕ}} — \emph{ārthika} — \textquotedblleft{}economic; financial\textquotedblright{}.

\begin{quote}
\kn{ಆ} \kn{ಕುಟುಂಬಕ್ಕೆ} \kn{ಆರ್ಥಿಕ} \kn{ತೊಂದರೆ} \kn{ಇದೆ}.
\end{quote}

\begin{quote}
\emph{ā kuṭumbakke ārthika tondare ide.}
\end{quote}

\begin{quote}
That family has money troubles.
\end{quote}

\kn{ಆರ್ಥಿಕತೆ}, the economy.

\begin{cousinweb}
\textbf{\kn{ಆರ್ಥಿಕ}} (\emph{ārthika}) is borrowed from Sanskrit \emph{ārthika} (Monier-Williams), an adjective from \emph{artha}, \textquotedblleft{}aim, wealth\textquotedblright{}.
\end{cousinweb}

\begin{grammarlens}[title={What you've built}]
One more word for giving reasons and opinions.
\end{grammarlens}

\subsection*{Your turn}
\begin{itemize}
\raggedright
  \item \emph{Say it:} \emph{ārthika}
  \item \emph{Say it:} \emph{ārthika}, once more
  \item \emph{Say it:} \emph{ā kuṭumbakke ārthika tondare ide.} — That family has money troubles.
  \item \emph{Recall:} say the Kannada for to waste, then the Kannada for development, then say \emph{ārthika} again
\end{itemize}

\begin{tcolorbox}[breakable,colback=teal!4,colframe=teal!35!black,title={Before you move on}]
What does \kn{ಆರ್ಥಿಕ} mean? (\textquotedblleft{}Economic; financial\textquotedblright{}.) Say it once more.
\end{tcolorbox}
\section[\texorpdfstring{\kn{ಹಣಕಾಸು} (haṇakāsu)}{ಹಣಕಾಸು (haṇakāsu)}]{\kn{ಹಣಕಾಸು} (haṇakāsu) — finance}

\label{lesson:KA-C401-hanakasu}



\begin{quote}
Before the new one: say the Kannada for development, then the Kannada for economic.
\end{quote}

\subsection*{You'll want to know: \kn{ಹಣಕಾಸು}}
\textbf{\kn{ಹಣಕಾಸು}} — \emph{haṇakāsu} — \textquotedblleft{}finance\textquotedblright{}.

\begin{quote}
\kn{ಮನೆಯ} \kn{ಹಣಕಾಸು} \kn{ಅಮ್ಮನೇ} \kn{ನೋಡಿಕೊಳ್ಳುತ್ತಾಳೆ}.
\end{quote}

\begin{quote}
\emph{maneya haṇakāsu ammanē nōḍikoḷḷuttāḷe.}
\end{quote}

\begin{quote}
Mother herself looks after the household finances.
\end{quote}

\kn{ಹಣ} (money) + \kn{ಕಾಸು} (coin).

\begin{cousinweb}
\textbf{\kn{ಹಣಕಾಸು}} (\emph{haṇakāsu}) is built from \emph{haṇa}, \textquotedblleft{}money\textquotedblright{}, paired with native \emph{kāsu}, \textquotedblleft{}a small coin, money\textquotedblright{} (DEDR 1431).
\end{cousinweb}

\begin{grammarlens}[title={What you've built}]
One more word for giving reasons and opinions.
\end{grammarlens}

\subsection*{Your turn}
\begin{itemize}
\raggedright
  \item \emph{Say it:} \emph{haṇakāsu}
  \item \emph{Say it:} \emph{haṇakāsu}, once more
  \item \emph{Say it:} \emph{maneya haṇakāsu ammanē nōḍikoḷḷuttāḷe.} — Mother herself looks after the household finances.
  \item \emph{Recall:} say the Kannada for development, then the Kannada for economic, then say \emph{haṇakāsu} again
\end{itemize}

\begin{tcolorbox}[breakable,colback=teal!4,colframe=teal!35!black,title={Before you move on}]
What does \kn{ಹಣಕಾಸು} mean? (\textquotedblleft{}Finance\textquotedblright{}.) Say it once more.
\end{tcolorbox}
\section[\texorpdfstring{\kn{ಉತ್ತೇಜಿಸು} (uttējisu)}{ಉತ್ತೇಜಿಸು (uttējisu)}]{\kn{ಉತ್ತೇಜಿಸು} (uttējisu) — to promote; to boost}

\label{lesson:KA-C401-uttejisu}



\begin{quote}
Before the new one: say the Kannada for economic, then the Kannada for finance.
\end{quote}

\subsection*{You'll want to know: \kn{ಉತ್ತೇಜಿಸು}}
\textbf{\kn{ಉತ್ತೇಜಿಸು}} — \emph{uttējisu} — \textquotedblleft{}to promote; to boost\textquotedblright{}.

\begin{quote}
\kn{ಸರ್ಕಾರ} \kn{ಸಣ್ಣ} \kn{ಉದ್ಯಮಗಳನ್ನು} \kn{ಉತ್ತೇಜಿಸಿತು}.
\end{quote}

\begin{quote}
\emph{sarkāra saṇṇa udyamagaḷannu uttējisitu.}
\end{quote}

\begin{quote}
The government boosted small businesses.
\end{quote}

Past (I …): \kn{ಉತ್ತೇಜಿಸಿದೆ}, \emph{uttējiside}.

\kn{ಉತ್ತೇಜನ}, encouragement.

\begin{cousinweb}
\textbf{\kn{ಉತ್ತೇಜಿಸು}} (\emph{uttējisu}) is built on Sanskrit \emph{uttejana}, \textquotedblleft{}incitement, encouragement, stimulation\textquotedblright{} (Monier-Williams), with the Kannada verb ending -\emph{isu}.
\end{cousinweb}

\begin{grammarlens}[title={What you've built}]
One more word for giving reasons and opinions.
\end{grammarlens}

\subsection*{Your turn}
\begin{itemize}
\raggedright
  \item \emph{Say it:} \emph{uttējisu}
  \item \emph{Say it:} \emph{uttējisu}, once more
  \item \emph{Say it:} \emph{sarkāra saṇṇa udyamagaḷannu uttējisitu.} — The government boosted small businesses.
  \item \emph{Recall:} say the Kannada for economic, then the Kannada for finance, then say \emph{uttējisu} again
\end{itemize}

\begin{tcolorbox}[breakable,colback=teal!4,colframe=teal!35!black,title={Before you move on}]
What does \kn{ಉತ್ತೇಜಿಸು} mean? (\textquotedblleft{}To promote; to boost\textquotedblright{}.) Say it once more.
\end{tcolorbox}
\section[\texorpdfstring{\kn{ಹೂಡು} (hūḍu)}{ಹೂಡು (hūḍu)}]{\kn{ಹೂಡು} (hūḍu) — to invest}

\label{lesson:KA-C401-hudu}



\begin{quote}
Before the new one: say the Kannada for finance, then the Kannada for to promote.
\end{quote}

\subsection*{You'll want to know: \kn{ಹೂಡು}}
\textbf{\kn{ಹೂಡು}} — \emph{hūḍu} — \textquotedblleft{}to invest\textquotedblright{}.

\begin{quote}
\kn{ನಾನು} \kn{ಉಳಿಸಿದ} \kn{ಹಣವನ್ನು} \kn{ಒಂದು} \kn{ಅಂಗಡಿಯಲ್ಲಿ} \kn{ಹೂಡಿದೆ}.
\end{quote}

\begin{quote}
\emph{nānu uḷisida haṇavannu ondu aṅgaḍiyalli hūḍide.}
\end{quote}

\begin{quote}
I invested the money I had saved in a shop.
\end{quote}

Past (I …): \kn{ಹೂಡಿದೆ}, \emph{hūḍide}.

\kn{ಹೂಡಿಕೆ}, an investment. Also to yoke: \kn{ಎತ್ತು} \kn{ಹೂಡು}.

\begin{cousinweb}
\textbf{\kn{ಹೂಡು}} (\emph{hūḍu}) is native Dravidian: DEDR 4361 lists Kannada \emph{pūḍu}, \textquotedblleft{}to join, to tie to, to yoke, to fix, to arrange\textquotedblright{}; old \emph{p-} became \emph{h-}, and \textquotedblleft{}to put money to work\textquotedblright{} is the modern sense.
\end{cousinweb}

\begin{grammarlens}[title={What you've built}]
One more word for giving reasons and opinions.
\end{grammarlens}

\subsection*{Your turn}
\begin{itemize}
\raggedright
  \item \emph{Say it:} \emph{hūḍu}
  \item \emph{Say it:} \emph{hūḍu}, once more
  \item \emph{Say it:} \emph{nānu uḷisida haṇavannu ondu aṅgaḍiyalli hūḍide.} — I invested the money I had saved in a shop.
  \item \emph{Recall:} say the Kannada for finance, then the Kannada for to promote, then say \emph{hūḍu} again
\end{itemize}

\begin{tcolorbox}[breakable,colback=teal!4,colframe=teal!35!black,title={Before you move on}]
What does \kn{ಹೂಡು} mean? (\textquotedblleft{}To invest\textquotedblright{}.) Say it once more.
\end{tcolorbox}
